% Chapter 45: Measurement, Probability, and So-Called Collapse (Complete Sample Chapter)
\chapter[Measurement, Probability, and Collapse]{Measurement, Probability, and So-Called Collapse}

Chapter~43 gave a prediction machine: state, amplitudes, squared moduli, outcome probabilities. Chapter~44 explained its between-measurement operation through continuous Schrödinger evolution. But what happens at measurement itself? Textbooks' intimidating collapse names the transition from many possibilities to one result. Is it a physical process or a bookkeeping rule? Must someone look? We open this black box and sort its contents into repeated experimental facts, mathematical calculation rules, and disputed interpretations. You may not leave knowing what collapse really is, but should recognize nine-tenths of its popular misreadings.

\step{A Counterintuitive Opening}

Light your phone screen and look through polarized sunglasses, common for driving and fishing. Rotate phone or lens until the screen blackens. Unsurprising: LCD light has aligned oscillation directions, and a polarizer resembles a fence passing oscillations along its gaps. Perpendicular fence and oscillation block everything.

Now the strange part. Keep the dark configuration fixed and find another polarizer, perhaps a lens removed from other sunglasses. Insert it about forty-five degrees between screen and first lens. Common sense says more filters only darken: two crossed fences already block everything, so how can an intervening diagonal fence brighten it?

Yet a clear bright patch returns, varying as the middle sheet rotates. Stranger still, move that sheet outside the pair, before or after, and nothing happens: darkness remains. Only the middle position works.

A filter blocks in some positions but apparently releases light in others. No defective lenses: a deep quantum secret lies between your screen and two pairs of sunglasses. By the end you will explain it fully and see its connection to knowledge, updating beliefs, and reality before measurement.

\step{Make a Guess First}

Record specific predictions before continuing:

Two crossed polarizers block everything. Insert a forty-five-degree third: (a) still black, (b) slightly brighter, or (c) clearly brighter? For (b) or (c), ask further: the third blocks some light. Was that light originally able to reach your eye, or already condemned by the other pair? If the latter, how can a filter resurrect nonexistent light?

Second, dim light until on average only one photon approaches the three sheets at a time. Indivisible, it passes whole or is absorbed whole. Does a photon facing two crossed sheets certainly die or have a half chance to survive? What changes with the middle sheet?

Finally, tomorrow's checkup includes a nurse measuring blood pressure. Does fastening the cuff change your pressure, and by how much? Apparently unrelated, but whether measurement disturbs its object becomes the threshold between everyday and quantum worlds later.

Keep your answers between the pages and compare afterward: which intuitions survive, which need rebuilding?

\step{Hands-on Experiment}

\begin{experiment}[Three Polarizers: a Filter That Releases Light]
\textbf{Materials}: an LCD phone or monitor with polarized output; polarized sunglasses; another polarizer, from other sunglasses, polarized 3D glasses, or the film before a discarded calculator LCD.

\textbf{Step one}: display plain white, for example an empty notes screen. Look through a polarizer and rotate to maximum darkness, showing perpendicular screen polarization and transmission direction.

\textbf{Step two}: keep it fixed, add and rotate a second sheet, and find darkness again. The two fences are now crossed, blocking all light.

\textbf{Step three}: keep both fixed and insert a third between them, rotating slowly. Brightness rises from black, peaking around forty-five degrees. Record the angle.

\textbf{Step four}: remove the third and place it before or after the pair. Every angle remains dark.

\textbf{Record and think}: why does the same sheet revive light only in the middle? Explain what it does to light and why a mere blocking filter would fail in the wrong position.
\end{experiment}

If three sheets are unavailable, two reproduce most observations. Even rotating one pair of sunglasses before a screen confirms its polarized output. Many longtime wearers never notice this daily quantum encounter.

\step{The Old Model Runs into Trouble}

Calculate with our best classical model, rope waves. Oscillation runs along the rope and fences pass gap-aligned components. A first vertical fence prepares vertical screen light; a second horizontal fence blocks it entirely. So far fine. A forty-five-degree third passes diagonal components, but now no components remain: the second already consumed the vertical oscillation. Refiltering zero remains zero. Rope waves decisively predict continued darkness. Experiment brightens: first failure.

Perhaps concede that the third rotates rather than filters? Classical waves can accommodate this, since a diagonal component includes a horizontal one. But single photons expose the real problem. Every photon facing two crossed sheets is absorbed, none detected; inserting the middle sheet lets a substantial fraction survive whole. There is no intermediate between absorption and arrival. The second condemned each photon, yet the third creates survivors. Filters cannot revive vanished things unless the first sheet's vertical preparation is no permanent identity and the middle does more than select: it rewrites the state.

Return to blood pressure. Cuff inflation indeed slightly changes blood flow, but less than the instrument's resolution; a thermometer removes negligible tongue heat. Classical physics assumes measurement reads existing properties, with disturbance arbitrarily reducible toward zero. Everyday success promotes this to an axiom: objects always possess definite position, velocity, and polarization simultaneously; measurement copies answers. Quantum mechanics founders here. If polarization were definite throughout, after the second sheet reads no horizontal component the third could not conjure one. The escape is admitting copying fails microscopically: measurement is a real physical interaction, replacing as well as revealing states.

A second wreck is clearer still. Chapter~43's individual electrons make double-slit fringes. Add detectors recording each route and fringes vanish into two humps. Perhaps a violent detector collision scrambles electrons? Make the observation gentle with faint light: lighter glimpses gradually weaken fringes rather than abruptly remove them, with information and visibility loss tracking exactly. The decisive factor is not collision force but whether which-path information exists. Nobody need read it: a record readable in principle destroys interference. Information itself is physical, beyond the old ledger.

Quantify individual photons. A one-milliwatt red laser emits about $3\times10^{15}$ photons per second, each around $3\times10^{-19}$~J. One per second on average requires fifteen orders of attenuation. Extraordinary sounding, yet laboratory attenuators and single-photon sources make individual double-slit and polarization experiments routine. Each photon fate above has an experimental receipt, not merely an imagined script.

\step{Inventing New Concepts}

Redesign the bankrupt ledger as physicists did, with four columns.

First, \textbf{observable}. Each measurement asks a chosen question: vertical polarization or forty-five-degree polarization? Same photon, different account, a crucial distinction. Everyday questionnaires hint at this: supporting reform versus tolerating its costs can yield opposite responses from the same people. But that is wording and psychology, not changing preexisting answers. Quantum questions instead choose physically different interactions.

Second, \textbf{possible outcomes}. Each question has a finite or infinite menu. Polarization's menu is short: pass or block, never half-pass. Individual outcomes are whole; probabilities emerge across repetitions.

Third, \textbf{probabilities}. The state, Chapter~43's wavefunction and amplitudes, plus Born gives each outcome's squared-amplitude probability. It promises proportions, not individual outcomes. Once resisted, this is now overwhelmingly tested: identically prepared repeated experiments follow Born statistics precisely.

Fourth, our protagonist, \textbf{state update}. After an outcome, future behavior follows a new compatible state, not the old one. A photon passing a forty-five-degree sheet must use that polarization for all later predictions; its old vertical information is canceled by the interaction. Replacing the state by the observed outcome's state is collapse's entire operational content. No column mentions anyone seeing. A detector current pulse, a blackened film grain, or an air molecule's altered path suffices when interaction leaves an indelible record. Consciousness is not even a footnote.

This resembles familiar Bayesian updating. A drawer contains one black sock and one white; you blindly take black, making the remaining sock's white probability jump from one-half to one. New evidence updates knowledge by deleting incompatible possibilities and renormalizing, formally like quantum updates. But mark two limits immediately. Sock color preexists; knowledge changes, not the sock. A forty-five-degree polarizer physically prepares a new photon state governing all later experiments, with no old-state photocopy available. Also, asking sock color then texture is order-independent; quantum vertical-then-horizontal and horizontal-then-vertical produce different accounts and distributions. Your middle-only polarizer already demonstrates order's physical role. A prewritten master answer table awaits Chapter~48's Bell inequalities. Here record that quantum predictions and such a table cannot coexist mathematically.

\step{The Model in One Sentence}

\begin{oneline}
Quantum measurement does not peek at prewritten answers. It physically interacts: you choose a question, nature supplies a Born-random answer, and the system enters a compatible new state from which all future predictions must restart.
\end{oneline}

Like shuffling, each measurement genuinely rearranges the deck. Unlike revealing a fixed facedown card. Nor is this an apology for crude instruments, a license for consciousness creating reality, or surrender to universal uncertainty. It supplies nature's most precise probabilistic predictions.

\step{The Mathematical Translator}

Translate the opening numerically. Let screen polarization be $0^\circ$. Chapter~43 decomposes states along orthogonal directions with amplitudes as coefficients. A $0^\circ$ photon, relative to a $45^\circ$ question, is
\[
|\,0^\circ\,\rangle=\cos 45^\circ\,|\,45^\circ\,\rangle+\sin 45^\circ\,|\,135^\circ\,\rangle .
\]
Four questions: it expresses the same state through different questions. Coefficients are geometric projections, cosine and sine, as with inclined-plane vector components in Chapter~4. It holds for ideal polarizers and monochromatic polarized light; real absorbing, scattering lenses make it an excellent approximation.

Born gives transmission through $45^\circ$ as $\cos^2 45^\circ=\tfrac12$. The transmitted state updates to $|\,45^\circ\,\rangle$. Now that $45^\circ$ photon meets $90^\circ$, with another $\cos^2 45^\circ=\tfrac12$. A screen photon's complete passage is
\[
P=\frac12\times\frac12=\frac14 .
\]
Check extremes: remove the middle sheet and $0^\circ$ meets $90^\circ$, giving $\cos^2 90^\circ=0$, darkness. Align it with the first and $\cos^2 0^\circ=1$ times $\cos^2 90^\circ=0$ remains dark. Only oblique placement brightens, matching observation; passing extremes earns confidence in the middle case.

The rule also describes smooth rotation. Incident polarization at $\theta$ to the transmission axis passes with $\cos^2\theta$ probability, hence bulk brightness proportional to $\cos^2\theta$. Malus's two-century empirical law becomes a direct Born consequence. Check $\theta=0^\circ$, all pass; $\theta=90^\circ$, all block; $\theta=45^\circ$, half pass. An old macroscopic law finds its single-photon origin, one of physics' rare pleasures.

\begin{figure}[htbp]
\centering
\begin{tikzpicture}[>=Stealth,scale=1.0]
  \draw[->,very thick] (-0.6,0) -- (8.2,0) node[right]{Light};
  \foreach \x/\a in {1.2/0, 3.8/45, 6.4/90}{
    \draw[thick,fill=gray!25] (\x,-0.85) rectangle ++(0.18,1.7);
    \node[below] at (\x+0.09,-0.9) {$\a^\circ$ sheet};
  }
  \node[above] at (-0.2,0.15) {\shortstack{Screen\\light}};
  \node[above] at (2.5,0.15) {$\tfrac12$ passes};
  \node[above] at (5.1,0.15) {$\tfrac12$ passes};
  \node[above] at (8.4,0.15) {$\tfrac14$ arrives};
\end{tikzpicture}
\caption{Three-polarizer quantum accounts: each sheet asks a question and updates transmitted states; sequential probabilities multiply. Removing the middle sheet drops total probability from $\tfrac14$ to $0$.}
\end{figure}

\begin{mathbridge}
\textbf{Conditional probability and renormalization.} Updating is mathematically almost boring: delete incompatible components and renormalize. Write $|\psi\rangle=a\,|\text{pass}\rangle+b\,|\text{block}\rangle$, with $|a|^2+|b|^2=1$. Observing pass leaves $|\text{pass}\rangle$: delete $b$, while the $a$ term is already unit length, sparing division. Sequential passage through $45^\circ$ then $90^\circ$ multiplies probabilities, ordinary conditional probability $P(A\,\text{and}\,B)=P(A)\,P(B\mid A)$. Quantum mechanics invents no new probability theory here, only \textbf{where probabilities come from}: amplitude moduli, not ignorance of hidden suits.

\textbf{Uncertainty relates widths.} Chapter~44 described position and momentum accounts as two versions of one state: narrower position packets spread momentum. Heisenberg expresses this as
\[
\Delta x\,\Delta p\geq \frac{\hbar}{2},\qquad \hbar\approx 1.05\times10^{-34}\ \text{J·s}.
\]
Read its subject carefully: $\Delta x$ and $\Delta p$ are not individual measurement errors, but distribution widths from position and momentum measurements on \textbf{identically prepared systems}. Perfect instruments cannot reduce their product below $\hbar/2$. A narrow-packet state itself carries no sharper momentum information. This is built into statehood, not engineering.
\end{mathbridge}

See how this rule separates scales. A ten-micrometer dust grain weighs about $5\times10^{-13}$~kg. Localize it to $\Delta x=10^{-6}$~m, roughly a hair's width:
\[
\Delta v\geq\frac{\hbar}{2m\,\Delta x}\approx\frac{1.05\times10^{-34}}{2\times5\times10^{-13}\times10^{-6}}\approx10^{-16}\ \text{m/s}.
\]
At that uncertain speed, a year's drift is smaller than a nucleus. Dust can have position and velocity with no visible quantum trace. An electron of $9.1\times10^{-31}$~kg localized within hydrogen's $\Delta x\approx10^{-10}$~m instead gives
\[
\Delta v\geq\frac{\hbar}{2m\,\Delta x}\approx6\times10^{5}\ \text{m/s},
\]
comparable to its characteristic $10^6$~m/s. The classical question of exact present position and directed speed has no sufficiently definite measurable answer. One inequality divided by vastly different masses explains both lifelong classical intuition and its failure inside atoms.

\begin{challenge}
\textbf{Why order has mathematical consequences.} Call $0^\circ$, $45^\circ$, and $90^\circ$ questions measurements $A$, $B$, and $C$. Consecutive $A$, $C$ always blocks; $A$, $B$, $C$ passes $\tfrac14$. Inserting $B$ changes the result: order matters. Chapter~46 translates measurements into operators, compressing this into noncommuting products, $AB\neq BA$. Their difference, the commutator, controls the uncertainty bound. Uncertainty and measurement order are two sides of one coin.

\textbf{Decoherence: the tireless environmental observer.} Why does dust never display a two-place superposition? Not abstract size but constant company: roughly $10^{18}$ to $10^{20}$ air-molecule collisions per second, plus reflected and emitted photons. Every collision effectively measures position, scattering phase information irretrievably into molecular motion. Classic Joos--Zeh estimates give a ten-micrometer grain in a one-centimeter-separated superposition only about $10^{-30}$ seconds before interference disappears, dozens of orders below instrumental resolution. The environment continuously measures with overwhelming efficiency, turning phase coherence into effectively classical probability mixtures. Classical laws have not taken over; dynamics makes appearances classical.
\end{challenge}

\step{The Model's Limits}

Our most important sorting exercise: collapse into facts, model, and interpretation.

\textbf{Facts}: repeated identical preparations obey Born; some sequences depend on order; individual outcomes are unpredictable while statistics are precise; measurement needs interaction, not human spectators, as tracks and photographic grains attest. A century of experiments has not shaken these.

\textbf{Model}: state vectors, Born probabilities, and projection updates deleting incompatible components and renormalizing. Imperfect or indirect measurements have generalized language, POVMs; here only ideal projective measurements were used. Where quantum mechanics applies these rules have never failed, but they do not explain why this particular result occurs or where interaction becomes measurement. That measurement problem remains a real theoretical gap.

\textbf{Interpretations}: different stories share everyday predictions, so distrust anyone selling one as experimentally proven. A neutral map: Copenhagen treats collapse as rule updating upon irreversible records, using a practical quantum/classical boundary without specifying it. Many worlds deletes collapse and keeps Schrödinger, with measurements branching the universe including you; branch probability is its weakness. Objective-collapse theories add small spontaneous stochastic collapses, localizing large bodies; these exceed interpretation because altered predictions face increasingly restrictive tests. Bohm gives particles definite positions guided by wavefunctions, at explicitly nonlocal cost, reproducing standard predictions. QBism makes states agents' beliefs about future experiences and collapse Bayesian updating, taking the sock analogy furthest while accepting order and no master answer table. Five stories, one prediction machine. Distinguishing common predictions from interpretation is your most valuable defense in quantum popularizations.

Three boundaries remain. Decoherence does not itself settle measurement: it explains lost interference and definite pointer directions, making superpositions look like classical alternatives, yet all options remain listed. Whether it alone explains one actual result remains disputed. Second, uncertainty belongs to states, not crude instruments. Third, collapse is no superluminal switch: one-particle measurement updates its state; Chapter~48 explains why entangled pairs still cannot communicate this way.

\begin{figure}[htbp]
\centering
\begin{tikzpicture}[>=Stealth,scale=1.0]
  % Before measurement
  \draw[->] (-0.2,0) -- (2.6,0);
  \draw[->] (0,-0.2) -- (0,1.8);
  \draw[fill=blue!40] (0.5,0) rectangle ++(0.6,1.2);
  \draw[fill=blue!40] (1.5,0) rectangle ++(0.6,1.2);
  \node[below] at (0.8,0) {Pass};
  \node[below] at (1.8,0) {Block};
  \node[above] at (1.3,1.9) {Before: each $\tfrac12$};
  % Arrow
  \draw[->,very thick] (3.0,0.8) -- (4.2,0.8) node[midway,above]{Pass observed};
  % After measurement
  \draw[->] (4.6,0) -- (7.4,0);
  \draw[->] (4.8,-0.2) -- (4.8,1.8);
  \draw[fill=orange!60] (5.3,0) rectangle ++(0.6,1.2);
  \draw[fill=gray!20] (6.3,0) rectangle ++(0.6,0.03);
  \node[below] at (5.6,0) {Pass};
  \node[below] at (6.6,0) {Block};
  \node[above] at (6.1,1.9) {After: $1$ and $0$};
\end{tikzpicture}
\caption{State update deletes incompatible possibilities and renormalizes after an outcome. Formally Bayesian, but physically the system itself changes state.}
\end{figure}

\step{Explaining the Opening Phenomenon}

Return to the light-releasing filter with complete accounts. Screen light starts at $0^\circ$. Sunglasses ask the $90^\circ$ question: $0^\circ$ and $90^\circ$ are exclusive, giving $\cos^2 90^\circ=0$. Every photon is absorbed and darkness results. No hidden horizontal component remains: this state's sole answer is block.

Insert $45^\circ$ between the questions. Asked $45^\circ$, the $0^\circ$ state is no longer certain: amplitude $\cos 45^\circ$ gives half passing. Crucially, they update to $45^\circ$, canceling the old file. The final sheet no longer receives doomed $0^\circ$ photons but ones with half-chance at $90^\circ$. Total $\tfrac14$ restores brightness. We added not another sieve but state-changing measurement. The filter resurrects no dead light; it changes the ledger defining death.

Wrong-position failure is equally clear: outside insertion leaves $0^\circ$ immediately followed by $90^\circ$, still blocking. Order is physics, a domestic noncommuting-measurement demonstration to return more sharply in Chapter~47's Stern--Gerlach experiment.

Quantum secure communication makes this engineering. Send individual polarization-encoded bits using two randomly chosen incompatible bases. An interceptor must choose a question, correct only half the time. The wrong measurement prepares a wrong state, and forwarding cannot erase its trace. Publicly comparing a small bit sample reveals about one-quarter errors. No sophisticated device principle is needed beyond these rules: measurement rewrites states and leaves statistical footprints. The eavesdropper loses to nature's ban on read-only quantum measurement, not inferior equipment.

\step{Thought Experiments and Challenges}

\begin{thought}[A Boxed Cat, Then the Moon Through Slits]
Schrödinger's cat was intended not to assert half-life/half-death but to X-ray measurement: if atom--detector--poison--cat amplifies microscopic superposition indefinitely, must macroscopic superposition follow? We have two responses. Dynamically, astronomical molecular and photon interactions decohere alive/dead phase information unimaginably fast, making it practically equivalent to classical alternatives, beyond any cat-outcome interferometer. Conceptually, classical appearance does not select one result, leaving the measurement gap open. More extreme, make the Moon interfere through slits: in principle merely ensure no photon, dust grain, or gravitational disturbance records its route. Decoherence estimates reveal the universe itself as a nonstop measuring instrument. The unwatched Moon is a joke because it never lacks things watching it.
\end{thought}

\begin{puzzles}
\item Split a $90^\circ$ turn among four sheets at $0^\circ,30^\circ,60^\circ,90^\circ$. What total transmission results? For $n$ sheets with $90^\circ/n$ steps, what limit follows, and what distinguishes filtering from measurement? \emph{Hint}: multiply each $\cos^2\theta$. For large $n$, $\cos^2(90^\circ/n)\approx1-(90^\circ/n)^2$, converting degrees to radians before multiplying. What does the exponential approach? Enough gentle questions block almost nothing.
\item A zero-disturbance microscope allegedly sees electrons without touching them, measuring position and momentum exactly to evade uncertainty. Refute it without formulas. \emph{Hint}: uncertainty concerns state widths, not disturbance. Measure position on half an identical ensemble and momentum on the rest; perfect instruments still yield a bounded width product. Evasion requires overturning states, not gentler looking.
\item Drawing a black sock raises the remaining sock's white probability from $\tfrac12$ to $1$, formally like a photon update. Give two substantive differences and identify a directly testable one. \emph{Hint}: sock colors preexist and order is irrelevant; quantum order changes distributions, as three polarizers demonstrate. Master prewritten answers await Chapter~48's Bell inequalities.
\item Why cannot an eavesdropper avoid detection with an extremely gentle measurement? \emph{Hint}: she must choose a basis, wrong half the time; wrong choice rewrites the state and forwarding cannot restore it. Sampling exposes roughly one-quarter mismatch. Gentleness cannot remove basis choice.
\end{puzzles}
